\subsection{End-to-End Solution and Recovery Workflow}
\label{app:native_solve_workflow}
The continuous objective certificate and the integral performance result solve
different tasks. We specify the finite branch family and original coordinates,
use working LPs and structured transformations to construct source points, and
certify their objectives against a complete-family bound. Hardware search and
schedule recovery produce a separate analytical execution. The following is
the implemented workflow, not a proposed complete branch-and-price algorithm.

\paragraph{1. Freeze the mathematical instance.}
The native family has 21 branches, 61,985,244 independent invocations,
62,030,334 logical nodes including conversions, and 68,907,890 storage
allocations. Its largest current branch has 15,044,386,833 implementation
choices before temporal expansion; the event domain has 5,809,955,157 labels.
These are logical domain counts, not an allocated sparse matrix.
Invocation membership, types, demand caps, hardware ranges, costs and row
recipes identify the LP. A generator may reuse a coefficient formula without
forcing its invocations to share decisions. The more general implementation
resolver can describe choices outside this frozen pool; those choices are
outside the present certificate.

\paragraph{2. Solve manageable working supports with OCS.}
Activated columns and temporal fields define a restricted working LP with a
checked lift to original coordinates. Let $f(x)=c^Tx+I_{[\ell,u]}(x)$,
let $Kx$ stack equality and inequality row activities, and let
$g(z)=I_{\{z_{\rm eq}=b_{\rm eq},\ z_{\rm ineq}\le b_{\rm ineq}\}}(z)$.
The selected fixed-step OCS kernel performs
\begin{align}
 \widehat v^{k+1}
   &=\operatorname{prox}_{\Sigma g^*}(v^k+\Sigma K\bar x^k),\nonumber\\
 \widehat x^{k+1}
   &=\operatorname{prox}_{T f}(x^k-TK^T\widehat v^{k+1}),\nonumber\\
 (x^{k+1},v^{k+1})
   &=\alpha_k(x^{\rm anchor},v^{\rm anchor})
       +(1-\alpha_k)(\widehat x^{k+1},\widehat v^{k+1}),
       \qquad \alpha_k=\frac1{k+2},\nonumber\\
 \bar x^{k+1}&=2x^{k+1}-x^k .
 \label{eq:native_halpern}
\end{align}
Each block restarts from its average. Centered arithmetic computes differences
from the anchor to avoid repeatedly summing a large baseline.
The primal proximal map is $\Pi_{[\ell,u]}(x-Tc)$; equality dual coordinates
receive the RHS shift, while $\le$-row multipliers additionally project onto
the nonnegative orthant. Reversible scaling and diagonal steps are verified
against the original operator. The tested paged adapter moves bounded vector
pages between host storage and GPU; it reduces GPU memory but still stores
the working iteration vectors somewhere.

\paragraph{3. Improve a full-domain primal without full vectors.}
The final continuation starts from an admitted implicit source point.
It retains identified parameter owners, changes legal original weighted
implementations, retains the embedding after its first real import, and
retires converted-host backups only after verifying their remaining consumers.
These are searches for better points of the same source LP, not new
constraints restricting all source solutions.
The accepted recipe translates a compact parent of $H$ events over $T$
delays, using weights $\omega_j=2(T-j)/(T(T+1))$.
Each linear state and lifecycle field is represented by its weighted sum;
duration is the maximum of the resulting resource-service fields.
The retained parent can exceed DRAM capacity, so convexity of feasible parents
is not the argument for feasibility. The checker proves the mixed capacity:
the largest mass of any $H$ consecutive delays is
$H(2T-H+1)/(T(T+1))$ for $H\le T$, with initial/static reservations accounted
for separately. Original readiness, boundary and service rows are also checked.
GPU blocks evaluate the two edges of the mixture; rational affine sums handle
the middle. An outward arithmetic allowance is included in the objective upper.

\paragraph{4. Certify the full LP objective.}
For $\min c^Tx$ subject to $Ex=e$, $Gx\le g$, $x\in[\ell,u]$,
any equality multipliers $\nu$ and legal inequality multipliers $\mu\ge0$ give
\[
 L=-e^T\nu-g^T\mu+
       \inf_{\ell\le x\le u}(c+E^T\nu+G^T\mu)^Tx .
\]
Box support and every residual sign are checked in original coordinates;
a small negative residual on an unbounded nonnegative axis cannot be rounded
away. A working-set dual is insufficient when omitted columns change this
infimum. The final certificate instead uses explicit coverage prices, work-mass
equalities and original service rows. Generator-wide inequalities establish
the required residual signs for all implementations; constant temporal
multipliers apply to every event. The minimum branch bound certifies the
perspective family. Appendix~\ref{app:native_lp_certificate} gives the actual
multipliers and endpoints. This provides approximate continuous objective
closure without completing a general arbitrary-multiplier pricing coordinator.

\paragraph{5. Search discrete hardware and recover software.}
The hardware branch-and-bound tree is separate from the whole-program branch
index $b$ in the LP. A tree node is a rectangle in a declared finite menu of
array counts, vector lanes, DRAM rates, host channels and local-service units.
Minimum-resource budget violations prune infeasible rectangles. Optimistic
resource values, relaxed private-capacity constraints and repriced local
implementations provide latency lower estimates. In a joint variant, each
parameter family has a two-state retained/streamed dynamic program. Nonnegative
prices on shared capacity rows yield the conditional Lagrangian bound
\[
 L(h,\mu)=\sum_f\min_{\pi\in\mathcal P_f}
       \left[t_f(h,\pi)+\sum_t\mu_t B_{ft}(\pi)\right]
       -\sum_t\mu_t C_t
\]
plus costs outside those paths. Here $f$ identifies an original parameter
family, $\pi$ an allowed retention path, and $B_{ft}$ its storage use.
GPU scans evaluate the independent paths; shared prices preserve the capacity
coupling in the bound. A path mixture or optimistic rectangle is not itself
an executable schedule.

The demonstrated menu contains 536,380,416 hardware combinations:
64--1,024 arrays in steps of two, 1,024--16,384 vector lanes in steps of 128,
512--2,560 GB/s DRAM in steps of 128 plus 2,608 GB/s, 1--64 host channels,
and 1--8 local-service units. An earlier fixed-software tree exhausted that
menu; the joint capacity-path tree made 1,190 GPU oracle calls and retained
142 unresolved schedule leaves, reduced to 78 by its subsequent replayed
incumbent. Its bounds were numerical and conditional on the then-current
cover, progress and implementation menu. They are not the final LP certificate,
and subsequent menu expansions require fresh tree bounds.

Only a replayed feasible design supplies a pruning incumbent; an unverified
candidate does not. Surviving hardware choices undergo integral cover/path
selection, capacity repair and transfer scheduling. Replay then computes
the predicted phase latencies using the selected chip's actual modeled
rates, clocks, formats and buffer reservations. This is how an integral
hardware/software result is obtained without claiming a completed tree over
every master selector.

\paragraph{What an exhaustive discrete proof would additionally require.}
One would need complete lower bounds at every remaining hardware/schedule node,
pricing over every allowed implementation under each branching decision,
and a fully explored or pruned frontier. Making tile or branch choices integral
alone does not make continuous rate/product envelopes exact. Spatial branching
to enforce physical products addresses a stronger nonlinear design problem;
it is not necessary to solve the frozen LP relaxation and is not claimed here.
Analytical replay supports the reported design latency, while fabrication and
global discrete optimality remain separate questions.
